\section{The game and the rules dialect studied}
\label{sec:rules}

Canadian Fish (also called Literature) is played by six players in two teams of
three, seated so that the teams alternate around the table. The deck is 54 cards,
partitioned into nine \emph{half-suits} of six cards each: the low band
$\{2,3,4,5,6,7\}$ and the high band
$\{9,10,\mathrm{J},\mathrm{Q},\mathrm{K},\mathrm{A}\}$ of each of the four suits,
together with a ninth half-suit containing the four eights and the two jokers.
The dealer is chosen at random, each player is dealt nine cards, and the player
to the dealer's left leads. A half-suit leaves play when it is claimed; the team
that takes at least five of the nine wins the game.

The dialect below is the one supplied by the project owner as the rules of
record, and the engine's \texttt{Rules} defaults match it on every point checked.
Published descriptions document mutually incompatible variants of several of
these choices \cite{pagat}; what follows is one explicit selection, not a
uniquely canonical rule set. It matches the strategy corpus this work draws on
\cite{develin}, and the contested choices are exposed as engine switches so that
their effect can be measured rather than assumed.

\paragraph{Asking.} On their turn a player asks one opponent for one named card.
The ask is legal only if the named card lies in a half-suit that is still live,
the asker holds at least one \emph{other} card of that half-suit, the asker does
not hold the named card, the target sits on the opposing team, and the target
still holds at least one card. If the target holds the named card they must
surrender it and the asker retains the turn; otherwise the turn passes to the
target. Every ask, its answer, every resulting transfer, and every player's card
count are public. The rules additionally license two public queries that the
engine models implicitly because both quantities are already in the public state:
the content of the previous two asks, and any player's remaining card count.

\paragraph{Declaring.} A claim on a half-suit is a \emph{declaration}: an exact
allocation naming, for each of the six cards, which member of the declaring
player's own team holds it. If every one of the six assignments is correct the
declaring team is awarded the half-suit; \emph{any} inaccuracy at all awards it
to the opposing team. Either way the six cards leave play, the half-suit becomes
inactive, and the turn returns to the player who held it before the declaration.
A declaration is legal at any moment, including during an opponent's turn. The
declarer need hold no card of the half-suit --- indeed need hold no cards at all
--- and may not declare on behalf of the opposing team.

\paragraph{Running out of cards.} A player with no cards leaves the asking cycle:
they can neither ask nor be asked, though they may still be named as the holder
of a card in a teammate's declaration and may still declare. If the turn reaches
a cardless player --- possible only after a declaration --- the turn passes to a
teammate who still holds cards. When an entire team is cardless, the other team
must declare every remaining half-suit with no further asking, and may misdeclare
and thereby hand half-suits to the cardless team.

\subsection{Two places where the engine is more restrictive than the rules}
\label{sec:rules-gap}

Both are coordination channels the rules permit and the engine does not use.
They are stated here because \S\ref{sec:mechanisms} reports that both were
measured and neither is worth building, which is a result about the game and not
only about the implementation.

\begin{enumerate}
  \item \textbf{Turn transfer.} The rules say that when a cardless player must
        pass the turn and both teammates hold cards, the team may openly
        strategise about which teammate receives it, but may share nothing beyond
        their willingness to receive it. The engine instead has the cardless
        player choose unilaterally from its own belief, and solicits no
        willingness bit. A genuine multi-candidate decision arises
        \numTransferRate{} times per game; \S\ref{sec:mech-rejected} reports what
        the channel is worth.
  \item \textbf{Declaration arbitration.} The rules do not say who acts first
        when two seats offer a declaration in the same instant. The engine
        resolves it by lowest seat, deliberately, to avoid comparing private
        confidences across seats. \S\ref{sec:mech-rejected} reports the measured
        cost of that choice.
\end{enumerate}

\subsection{Engine settings, and the switches that change them}
\label{sec:rules-dialect}

\begin{itemize}
  \item \textbf{Deck:} 54 cards, nine half-suits, nine per player
        (\texttt{-{}-sets=8} gives the 48-card, eight-half-suit form).
  \item \textbf{Declaration timing:} any seat may declare at any moment, in or
        out of turn, holding cards or not (\texttt{-{}-no-out-of-turn},
        \texttt{-{}-no-cardless-declare}).
  \item \textbf{Misdeclaration:} the half-suit is awarded to the opposing team;
        published variants differ here \cite{pagat}.
  \item \textbf{Simultaneous declarations:} the lowest seat index that offers a
        declaration acts first --- a modelling choice, not a rule
        (\texttt{-{}-arb=low\textbar high\textbar turn}).
  \item \textbf{Forced endgame:} an eight-level willingness ladder selects the
        declarer when one team is cardless (\texttt{-{}-legacy} gives the
        two-level v0.3 ladder). \S\ref{sec:diag-forced} reports that in the
        shipped policy this ladder is inert.
  \item \textbf{Action cap:} 400 asks, any residue adjudicated neutrally by
        majority physical holding (\texttt{-{}-maxasks}).
\end{itemize}

No prearranged signalling conventions are available to any policy in the
configuration measured here, and team communication is limited to the willingness
bit above and the successor choice made by a cardless turn-holder. The literature
records that pre-agreed conventions are explicitly forbidden in this game
\cite{develin} while also documenting a named convention in real use
\cite{pagat}, and a pair of policies trained together has an implicit
pre-agreement by construction; \S\ref{sec:mech-conventions} states how \vfive{}
handles that tension.

\paragraph{The action cap is a backstop, not a rule.} Two sufficiently patient
policies can stall indefinitely, and this report is largely about a configuration
that does. Its incidence is therefore reported with every experiment rather than
assumed to be zero: it is \numPatientCap\% for the unforced mirror configuration
of \S\ref{sec:diag-dilemma} and zero for the shipped \vfour{} configuration,
which avoids it by forcing declarations that \S\ref{sec:diag-dilemma} shows are
\numPostHorizonWrong\% wrong.
